%% LyX 2.4.4 created this file.  For more info, see https://www.lyx.org/.
%% Do not edit unless you really know what you are doing.
\documentclass[english]{article}
\usepackage[T1]{fontenc}
\usepackage[latin9]{inputenc}
\usepackage{enumitem}
\usepackage{amsmath}
\usepackage{amsthm}
\usepackage{amssymb}
\usepackage{cancel}
\usepackage{geometry}
\geometry{verbose,tmargin=1in,bmargin=1in,lmargin=1in,rmargin=1in}
\PassOptionsToPackage{normalem}{ulem}
\usepackage{ulem}

\makeatletter
%%%%%%%%%%%%%%%%%%%%%%%%%%%%%% Textclass specific LaTeX commands.
\numberwithin{equation}{section}
\numberwithin{figure}{section}
\newlength{\lyxlabelwidth}      % auxiliary length 

%%%%%%%%%%%%%%%%%%%%%%%%%%%%%% User specified LaTeX commands.
\usepackage{fancyhdr}
\usepackage{lastpage}
\pagestyle{fancy}
\usepackage{enumitem}
\usepackage{ifthen}
\everymath=\expandafter{\the\everymath\displaystyle}
%\DeclareMathSizes{10}{10}{10}{10}
\usepackage{multicol}

\usepackage{tikz}
\usetikzlibrary{patterns}
\usetikzlibrary{plotmarks}
\usepackage{pgfplots}
\pgfplotsset{compat=newest}

\usepackage{calc}
\usepackage[nomessages]{fp}

\usepackage{datetime}

%%%%%commands
\newcommand{\mb}[1]{\mathbb{#1}}
\newcommand{\funct}[3]{$#1\colon#2\rightarrow#3$}
% Added by lyx2lyx
\setlength{\parskip}{\smallskipamount}
\setlength{\parindent}{0pt}

\makeatother

\usepackage{babel}
\begin{document}
\renewcommand{\author}{Dr.\,James Ashe}

\newcommand{\classNum}{439}

\newcommand{\classSec}{A}

\newcommand{\nameType}{Solutions\hspace{-4cm}}

\newcommand{\examType}{Midterm}

\newcommand{\Date}{\formatdate{4}{3}{2014}} %%\AdvanceDate[12] \today

%\newdate{my_date}{\day}{\month}{\year}%   
%\AdvanceDate[-5] 
%\displaydate{my_date}

%%First page's header%%%%%%%%%%%%%%%%%%%%%%
\fancypagestyle{firststyle} %%header settings for first pagr
{
	\fancyhead[L]{MTH \classNum{}(\classSec{}) \author{}\\
	\textbf{\examType{}} \Date{}}
	\fancyhead[C]{\textbf{\nameType}}
	\setlength{\headsep}{14pt}
}
%%%%%%%%%%%%%%%%%%%%%%%%%%%%%%%%%%%%%%%%%%

%%Next page's header%%%%%%%%%%%%%%%%%%%%%%%
%\setlist{nolistsep} %eliminates space between list items
\lhead{MTH \classNum{}(\classSec{})} 
\chead{\textbf{\nameType}} 
\cfoot{\thepage{}~of~\pageref{LastPage}} 
\setlength{\headsep}{5pt} 
%%%%%%%%%%%%%%%%%%%%%%%%%%%%%%%%%%%%%%%%%%%

%%Various settings; see the preamble for other settings%%%%%
%%\setenumerate[0]{leftmargin=12pt,itemindent=0pt,labelwidth=0pt}
\renewcommand{\labelenumi}{\textbf{\arabic{enumi}.}}
\renewcommand{\labelenumii}{\textbf{\alph{enumii}.}}
\thispagestyle{firststyle}
%%%%%%%%%%%%%%%%%%%%%%%%%%%%%%%%%%%%%%%%%%
\newcommand{\xspace}{\vspace{.85cm}}

\small\emph{Unless stated otherwise, all rings are assumed to have
the normal addition and multiplication operations.}

\subsection*{Part 1}
\begin{enumerate}[resume]
\item {[}3 pts each{]} Provide the definition for each of the following:

\begin{enumerate}
\item Integral domain
\item Onto function
\item Ring isomorphism
\item Ideal

\begin{enumerate}
\item []See the notes.\xspace
\end{enumerate}
\end{enumerate}
\item {[}3 pts each{]} Provide an example for each of the following (do
not use an example from problems \ref{enu:first}-\ref{enu:last}):

\begin{enumerate}
\item A ring without identity

\begin{enumerate}
\item []The ring of even integers does not have a multiplicative identity
since $1$ is not even. Subrings of rings with identity provide examples
such as $n\mathbb{Z}$ for $n>1$.
\end{enumerate}
\item A non-commutative ring 

\begin{enumerate}
\item []The ring of $2\times2$ matrices with real entries (or with integer,
integer $\bmod n$, rational, or complex entries).
\end{enumerate}
\item A field 

\begin{enumerate}
\item []A few of the obvious ones: $\mathbb{R}$, $\mathbb{Q}$, $\mathbb{C}$.
\end{enumerate}
\item A function which is not 1-1.

\begin{enumerate}
\item []$f(x)=1$ (or any constant function) and $f(x)=x^{2}$ for $f\colon\mathbb{R}\rightarrow\mathbb{\mathbb{R}}$;
or $f(a)=a$ and $f(a)=b$ for $f\colon\{a\}\rightarrow\{a,b\}$. 
\end{enumerate}
\item An onto function

\begin{enumerate}
\item []$f(x)=x$ (the identity function) and $f(x)=x^{3}$ for $f\colon\mathbb{R}\rightarrow\mathbb{\mathbb{R}}$;
or $f(a)=a$ for $f\colon\{a\}\rightarrow\{a\}$.
\end{enumerate}
\item Two rings which are isomorphic.

\begin{enumerate}
\item []$\mathbb{R\times\mathbb{R}\cong\mathbb{C}}$, $\mathbb{Z}_{6}\cong\mathbb{Z}_{2}\times\mathbb{Z}_{3}$,
and of course any ring $R$ is isomorphic to itself, $R\cong R$.
\end{enumerate}
\item Two rings which are not isomorphic.

\begin{enumerate}
\item []$\mathbb{R}\cancel{\cong}\mathbb{C}$, $\mathbb{R}\cancel{\cong}\mathbb{Z}$,
$\mathbb{Q}\cancel{\cong}\mathbb{Z}$, $\mathbb{Z}_{3}\cancel{\cong}\mathbb{Z}_{4}$,
etc. It is easy to see that any two rings with different order are
not isomorphic.
\end{enumerate}
\item An ideal.

\begin{enumerate}
\item []$2\mathbb{Z}$ or $n\mathbb{Z}$ for any $n>1$; for and ring $R$,
$0_{R}$ and $R$ are ideals.
\end{enumerate}
\end{enumerate}
\end{enumerate}

\subsection*{\newpage}

\subsection*{Part 2}

\emph{{[}9 pts each{]} Choose the indicated number of problems from
each section.}

\subsubsection*{Rings: choose any three.}
\begin{enumerate}
\item \label{enu:first}Complete the following:

\begin{enumerate}[resume]
\item Give an example of a subset of the integers which \uline{is not} a
ring (under the usual addition and multiplication).

\begin{enumerate}[resume]
\item []The set of odd integers, $\mathbb{Z}-\{0\}$, $\{1,2,3,4\}$, etc.
\end{enumerate}
\item Give an example of a subset of the integers which \uline{is} a ring.

\begin{enumerate}[resume]
\item []$2\mathbb{Z}$ (the set of even integers), $n\mathbb{Z}$, $\{0\}$,
$\mathbb{Z}$ (I didn't specify that it had to be a proper subset),
etc.
\end{enumerate}
\end{enumerate}
\end{enumerate}
\begin{enumerate}[resume]
\item Prove that $\mathbb{Q}$ is subring of $\mathbb{R}$.

\begin{proof}
$\mathbb{Q}\subseteq\mathbb{R}$ and $\mathbb{Q}\neq\emptyset$ since
$0\in\mathbb{Q}$ (and many others). Let $\frac{a}{b},\frac{c}{d}\in\mathbb{Q}$
(so $a,b,c,d\in\mathbb{Z}$). $\frac{a}{b}-\frac{c}{d}=\frac{ad-bc}{bd}\in\mathbb{Q}$
and $\frac{a}{b}\frac{c}{d}=\frac{ac}{bd}\in\mathbb{Q}$. So $\mathbb{Q}$
is a subring by the subring test.
\end{proof}
\item Complete the following:

\begin{enumerate}[resume]
\item Prove that the set of odd integers is not a subring of $\mathbb{Z}$.

\begin{proof}
$3$ and $5$ are odd integers, but $3+5=8=2\cdot4$ is even. So the
set of odd integers is not a subring of $\mathbb{Z}$ since it is
not closed under addition.
\end{proof}
\item Prove that in general the union of any two rings is not a ring.

\begin{proof}
This is equivalent to showing that there are two rings such that their
union is not a ring. $\{0,1\}$ and $3\mathbb{Z}$ are subrings of
$\mathbb{Z}$, but $1+3=4\notin\{0,1\}\cup3\mathbb{Z}$. $M_{2\times2}(\mathbb{R})$
and $M_{3\times3}(\mathbb{R})$ are both rings, but their union certainly
is not -addition an multiplication between their elements isn't even
defined.
\end{proof}
\end{enumerate}
\item Prove that the intersection of any two subrings of a ring $R$ is
also a subring of $R$.

\begin{proof}
Let $A$ and $B$ be rings. Since $0_{R}\in A\cap B$; $A\cap B\neq\emptyset$
and we can choose $x,y\in A\cap B$. Since $x,y\in A$ and $A$ is
a ring, we have that $x-y\in A$ and $xy\in A$. Similarly $x-y\in B$
and $xy\in B$. Thus $x-y,xy\in A\cap B$, and $A\cap B$ is a subring
by the subring test. 
\end{proof}
\item Let $R$ be a ring and $a,b\in R$. Prove that $(-a)(-b)=ab$. Hint:
$(a+(-a))(-b)=0$.

\begin{proof}
Let $a,b\in R$. Using the hint, $(a+(-a))(-b)=0\Rightarrow a(-b)+(-a)(-b)=0$.
So $(-a)(-b)$ is the unique additive inverse of $a(-b)$ which is
$ab$. (Since $a(-b)+ab=a(-b+b)=0$).
\end{proof}

\begin{proof}
$(-a)(-b)=(-a)(-b)+(-a)b+(ab)=(-a)(-b+b)+(ab)=ab$.
\end{proof}
\end{enumerate}
\vfill

\subsubsection*{Integral Domains and Fields: choose any two.}
\begin{enumerate}[resume]
\item Prove that $\mathbb{Z}_{n}$ is \uline{not} an integral domain if
$n$ is not prime.

\begin{proof}
Since $n$ is not prime, $n=mq$ for two non-zero integers such that
$1<m,q<n$. $m\bmod n$ and $q\bmod n$ are not zero, but $mq\bmod n=n\bmod n=0\bmod n$.
\end{proof}
\item Prove that $\mathbb{Z}_{3}$ is a field. 

\begin{proof}
$3$ is prime so then $\mathbb{Z}_{3}$ is a field.
\end{proof}

\begin{proof}
We know that is $\mathbb{Z}_{3}$ is a ring. $1\cdot1=1$ and $2\cdot2=1\bmod3$.
So every element has a multiplicative inverse and $\mathbb{Z}_{3}$
must be a ring.
\end{proof}
\item Let $R$ be an integral domain and suppose that $ax-bx=0$. Prove
that if $a\neq b$ then $x=0$.

\begin{proof}
If $ax-bx=0$ then $(a-b)x=0$. Since $a\neq b$, $a-b\neq0$. So
$x=0$ since $R$ is an integral domain.
\end{proof}
\item Let $R$ be an integral domain and let $a,b,c\in R$ where $c\neq0$.
If $ac=bc$ then $a=b$.

\begin{proof}
$ac=bc\Rightarrow ac-bc=0\Rightarrow(a-b)c=0.$ Because $R$ is an
integral domain and $c\neq0$, $a-b$ must be zero. This implies that
$a-b$.
\end{proof}
\item Let $R$ be a field. Prove that $R$ is an integral domain.

\begin{proof}
Suppose that $ab=0$ and $a\neq0$ for some $a,b\in R$. Since $R$
is a field, there exists an $a^{-1}$ such that $a^{-1}ab=a^{-1}0\Rightarrow b=0$.
So then $R$ is an integral domain.
\end{proof}
\end{enumerate}
\vfill

\subsubsection*{Homomorphisms, isomorphisms, and ideals: choose any three.}
\begin{enumerate}[resume]
\item Prove that $f\colon\mathbb{R}\rightarrow\mathbb{R}$ where $f(x)=e^{x}$
is one-to-one but is not onto.

\begin{proof}
(1-1) Let $a,b\in\mathbb{R},$ $e^{a}=e^{b}\Rightarrow\ln(e^{a})=\ln(e^{b})\Rightarrow a=b$.
\\
(not onto) $-1\in\mathbb{R}$ but there is no real number such that
$e^{x}=-1$. 
\end{proof}
\item Let \funct{f}{R}{S} and \funct{g}{S}{T} for rings $R$, $S$, and
$T$. Prove that \funct{f\circ g}{R}{T} is a homomorphism.

\begin{proof}
$f(g(a+b))=f(g(a)+g(b))=f(g(a))+f(g(b))$, and $f(g(ab))=f(g(a)g(b))=f(g(a))f(g(b))$
since $g$ and $f$ are each homomorphisms.
\end{proof}
\item Let $f\colon\mathbb{Z}_{4}\rightarrow\mathbb{Z}_{4}$. 

\begin{enumerate}[resume]
\item If $f(x)=2x$ prove that $f$ is not a homomorphism.
\item If $f(x)=3x$ prove that $f$ is a homomorphism.
\end{enumerate}
\item Prove that $f\colon\sqrt{2}\mathbb{Z}\rightarrow\sqrt{2}\mathbb{Z}$
where $f(a+b\sqrt{2})=a-b\sqrt{2}$ is an isomorphism ($\sqrt{2}\mathbb{Z}=\{a+b\sqrt{2}\colon a,b\in\mathbb{Z}\}$).
\item \label{enu:last}Show that any ring isomorphism \funct{f}{\mb{Z}}{\mb{Z}}
must be the identity map.
\item Let $A$ and $B$ be ideals of $R$. Prove that $A\cap B$ is an ideal
of $R$.
\item Prove that $\mathbb{R}\times\mathbb{R}=\{(a,b)\colon a,b\in\mathbb{R}\}$
is isomorphic to $\mathbb{C}=\{a+b\sqrt{-1}\colon a,b\in\mathbb{R}\}$
if multiplication in $\mathbb{R}\times\mathbb{R}$ is defined as $(a,b)(c,d)=(ac-bd,ad+bc)$
and addition is defined as $(a,b)+(c,d)=(a+c,b+d)$.
\end{enumerate}
\vfill
\end{document}
